% Variáveis utilizadas na quantificação de produção acadêmica

% Nomes dos periódicos
\newcommand{\aps}{Asian Population Studies}
\newcommand{\csp}{Canadian Studies in Population}
\newcommand{\demography}{Demography}
\newcommand{\dr}{Demographic Research}
\newcommand{\ejp}{European Journal of Population}
\newcommand{\ijps}{International Journal of Population Studies}
\newcommand{\japs}{Journal of African Population Studies}
\newcommand{\pdr}{Population and Development Review}
\newcommand{\prpr}{Population Research and Policy Review}
\newcommand{\ps}{Population Studies}
\newcommand{\psp}{Population, Space and Place}
\newcommand{\rebep}{Revista Brasileira de Estudos de População}
\newcommand{\relap}{Revista Latinoamericana de Población}

% Quantificação da escassez de publicações de deficiência, surdez e língua de sinais
\newcommand{\wosSUMoverall}{16.575}
\newcommand{\wosSUMdisab}{148}
\newcommand{\wosSUMdeaf}{4}
\newcommand{\wosSUMsign}{2}

\newcommand{\wosRATIOdisab}{111}
\newcommand{\wosRATIOdeaf}{4143}
\newcommand{\wosRATIOsign}{8286}
